% =============================================================================
% CHAPTER 5: Conclusions and Future Work
% =============================================================================
% SPDX-FileCopyrightText: 2026 Frank Langbein <frank@langbein.org>
% SPDX-License-Identifier: LPPL-1.3c
%
% Suggested sections: Conclusions; Future work. Do not introduce new material in
% Conclusions---summarise aims and main results. Future work: next steps, open
% questions, what you would do with more time. Optional: short summary of
% contributions (as part of Conclusions or a bullet list).
% =============================================================================


\section{Conclusions}\label{sec:conclusions}

This project set out to investigate whether generative AI can be used to support practical cyber security education by creating vulnerable environments on demand. The motivation was that many CTF and cyber range platforms rely on static, pre authored challenges. These are valuable for structured learning, but their replayability can decrease once solutions are known, shared, or memorised. VulnForge explored a different approach: using an LLM to generate structured vulnerability specifications, deploying them into isolated Linux containers, allowing learners to interact with them over SSH, supporting them through a contextual tutor, validating flag submission, and cleaning up the environment afterwards.

The main achievement of the project is a working prototype that demonstrates this lifecycle, from producing structured vulnerabilities, to deploying them to user machines and verifying the exploit path. This shows that the project did not only produce isolated components, but combined them into a complete platform where generated challenges can be created, attempted and destroyed.

The evaluation suggests that this architecture is technically feasible, but not yet consistently reliable. In the full lifecycle evaluation, 57 out of 75 generated challenges completed successfully end to end. Generation, deployment, injection, and cleanup were generally strong, with cleanup succeeding for every container that reached that stage. This is an important result because the platform deliberately creates vulnerable systems, so failed or expired challenges must not be left running unnecessarily. The weaker results came from verification and exploit reliability. Some generated challenges could be deployed and appear healthy, but the generated environment, validation logic, and intended exploit path did not always align well enough for the challenge to be verified successfully.

This means the project supports the central claim that AI-assisted vulnerable environment generation can work in a controlled educational setting, but also shows why the LLM cannot be trusted as an unconstrained challenge author. A generated specification may be syntactically valid while still producing an unsuitable, unsafe, or unexploitable challenge. For that reason, the most important design lesson from VulnForge is that the AI component must be surrounded by guardrails, structured schemas, server side validation, policy checks, exploit checks, an cleanup mechanisms. The system is therefore not simply an LLM prompt wrapped in a web interface, but a constrained pipeline for turning uncertain generated content into testable training environments.

The project also showed that challenge reliability is affected by more than the nominal vulnerability category or difficulty level. SQL injection was the strongest category in the evaluation, while insecure file permissions was the weakest, partly because distinguishing intended exposure from accidental flag leakage is more difficult in that category. The difficulty results also did not show a simple pattern where harder challenges failed more often. Thi suggests that VulnForge can generate across different categories and levels, but it does not yet prove that those difficulty labels are pedagogically calibrated. Demonstrating that beginner, intermediate, and advanced challenges are experienced that way by learners would require a separate user centred evaluation.

Overall, VulnForge achieved its main aim as a feasibility prototype. It demonstrates that generated vulnerable environments can be integrated into a practical training workflow, and it provides evidence that such environments can be deployed, used, verified, and cleaned up automatically in many cases. At the same time, the evaluation identifies clear limitations in reliability, validation and infrastructure capacity. The project therefore shows both the promise and the risks of using generative AI for cyber security training. It can increase variety and reduce reliance on manually authored challenges, but only if the generated content is treated as untrusted and subjected to careful technical checks.
